\documentclass[../../../include/open-logic-section]{subfiles}

\begin{document}

\olfileid{sth}{replacement}{limofsize}
\olsection{કદની મર્યાદા}

પ્રતિસ્થાપનને ``આંતરિક'' રીતે વાજબી ઠરાવવાનો કદાચ સૌથી
પ્રચલિત પ્રયાસ નીચેના ખ્યાલથી શરૂ થાય છે:
\begin{enumerate}
	\item[] \limofsize. કોઈ પણ વસ્તુઓનો ગણ બને છે,
	જો તેમની સંખ્યા અતિશય મોટી ન હોય.
\end{enumerate}
આ સિદ્ધાંતથી પ્રતિસ્થાપન તરત જ યોગ્ય ઠરશે. છેવટે,
પ્રતિસ્થાપનથી બનતો ગણ, જે ગણમાંથી બન્યો હોય તેના કરતાં
વધુ મોટો હોઈ ન શકે. ચોક્કસ રીતે: ધારો કે પ્રતિસ્થાપનથી
$\funimage{\tau}{A} = \Setabs{\tau(x)}{x \in A}$ ગણ
બનાવો; તો $\cardle{\funimage{\tau}{A}}{A}$; તેથી $A$ના
!!{element}s ગણ રચવા માટે અતિશય વધુ ન હોય તો તેમની
પ્રતિમાઓની સંખ્યા પણ $\funimage{\tau}{A}$ ગણ રચવા
માટે અતિશય વધુ નહીં હોય.

પ્રતિસ્થાપનને વાજબી ઠરાવવા \limofsize{}નો આધાર
લેવામાં સ્પષ્ટ મુશ્કેલી એ છે કે આપણે હજી સુધી
\limofsize{} જેવો કોઈ સિદ્ધાંત \emph{જણાવ્યો નથી}.
વધુમાં, \crefrange{sth:story::chap}{sth:z::chap}માં
ગણની સંચયી-પુનરાવર્તી કલ્પના વિશે વાત કરી ત્યારે
\limofsize{} તરફનો કોઈ \emph{સંકેત પણ} મળ્યો નહોતો.
ખરેખર, આ જ કારણથી બૂલોઝે એક વાર લખ્યું હતું:
``કદાચ તારણ કાઢી શકાય કે ગણસિદ્ધાંતની `પાછળ' ઓછામાં
ઓછા બે વિચારો છે'' \citeyearpar[p.~19]{Boolos1989}.
એક તરફ, ગણની સંચયી-પુનરાવર્તી કલ્પનાની આસપાસના
વિચારો $\Z$ને વાજબી ઠરાવવા માટે છે. બીજી તરફ,
\limofsize{} પ્રતિસ્થાપનને વાજબી ઠરાવવા માટે છે.

પરંતુ પ્રશ્ન માત્ર એટલો નથી કે આપણે અત્યાર સુધી
\limofsize{} વિશે \emph{કંઈ કહ્યું નથી}. પ્રશ્ન તો
એ છે કે \limofsize{} (અહીં રજૂ કર્યું છે એ સ્વરૂપે)
ગણની સંચયી-પુનરાવર્તી કલ્પના સાથે બહુ ખરાબ રીતે
બંધબેસે છે. છેવટે, તેમાં ગણો \emph{તબક્કાઓમાં}
રચાય છે એ વિચારનો કોઈ ઉલ્લેખ નથી.

આ ``બે વિચારો''નો, એટલે ગણની સંચયી-પુનરાવર્તી
કલ્પના અને \limofsize{}નો, ઇતિહાસ જોતાં આમાં
બહુ આશ્ચર્ય નથી. અંતે આ ``બે વિચારો'' ગણસિદ્ધાંતના
વિરોધાભાસો રોકવાના બે સાવ જુદા પ્રયાસો છે.
ગણની સંચયી-પુનરાવર્તી કલ્પના રસેલનો વિરોધાભાસ
આશરે આ રીતે રોકે છે: \emph{રસેલનો ગણ અસ્તિત્વમાં હશે
એવી અપેક્ષા રાખવી જ ન જોઈએ, કારણ કે તે કોઈ પણ
તબક્કે ``રચાતો'' નથી}. બીજી તરફ, \limofsize{}
રસેલના ગણને આશરે આ રીતે નકારી કાઢે છે:
\emph{રસેલનો ગણ અસ્તિત્વમાં હશે એવી અપેક્ષા રાખવી જ
ન જોઈએ, કારણ કે તે અતિશય મોટો હશે}.

આ રીતે મૂકીએ તો સ્પષ્ટ કહીએ: વિરોધાભાસોના જવાબ તરીકે
\limofsize{}ને \emph{ઘણું} વધુ ઔચિત્ય જોઈએ. ઉદાહરણ
તરીકે, રસેલના વિરોધાભાસનું આ સ્વરૂપ વિચારો:
\emph{એવો કોઈ પગ જાતિનો કૂતરો નથી જે પોતાને ન સૂંઘતા
હોય એવા અને માત્ર એવા પગ જાતિના કૂતરાઓને સૂંઘે}
(જુઓ \olref[sth][story][rus]{sec}). જો તમે પૂછો કે
``એવો કૂતરો શા માટે નથી?'', તો એવો જવાબ સારો નહીં ગણાય
કે તેને અતિશય વધુ કૂતરાઓને સૂંઘવા પડે. તો રસેલનો
ગણ અસ્તિત્વમાં ન હોવાનું સહજ ખુલાસો એવો કેમ ગણાય કે
તે અસ્તિત્વમાં રહી શકે એટલા કરતાં ``અતિશય મોટો'' હશે?

ટૂંકમાં, \limofsize{}ની ``સહજ'' પ્રેરણા અંગે તમને
થોડી મૂંઝવણ થાય, તો તે સમજી શકાય તેવી છે.

\end{document}
